% Fórmulas ancla — Probabilidad y Estadística
% Docente Mgr Alfredo Delgadillo Cossio
% Fragmento compilable. Requiere amsmath.
\documentclass{article}
\usepackage[T1]{fontenc}
\usepackage[utf8]{inputenc}
\usepackage[spanish]{babel}
\usepackage{amsmath,amssymb}
\title{Fórmulas ancla --- Probabilidad y Estadística}
\author{Docente Mgr Alfredo Delgadillo Cossio}
\date{}
\begin{document}
\maketitle

\section{Unidad 1}
\[
f_i = \frac{n_i}{n},\qquad
x_i^* = \frac{L_i+U_i}{2},\qquad
k \approx 1 + 3.322\log_{10} n
\]

\section{Unidad 2}
\begin{align*}
\bar{x} &= \frac{1}{n}\sum_{i=1}^n x_i, &
G &= \exp\Bigl(\frac{1}{n}\sum \ln x_i\Bigr), &
H &= \frac{n}{\sum (1/x_i)} \\
s^2 &= \frac{1}{n-1}\sum (x_i-\bar{x})^2, &
\mathrm{CV} &= \frac{s}{\bar{x}}, &
z_i &= \frac{x_i-\bar{x}}{s}
\end{align*}

\section{Unidad 3}
\[
\hat{y}=a+bx,\quad
b=\frac{\sum(x_i-\bar{x})(y_i-\bar{y})}{\sum(x_i-\bar{x})^2},\quad
r=\frac{\sum(x_i-\bar{x})(y_i-\bar{y})}{\sqrt{\sum(x_i-\bar{x})^2\sum(y_i-\bar{y})^2}}
\]

\section{Unidad 4}
\[
P(A\mid B)=\frac{P(A\cap B)}{P(B)},\qquad
P(A_i\mid B)=\frac{P(B\mid A_i)P(A_i)}{\sum_j P(B\mid A_j)P(A_j)}
\]
\[
P(n,k)=\frac{n!}{(n-k)!},\qquad
\binom{n}{k}=\frac{n!}{k!(n-k)!}
\]

\section{Unidad 5}
\[
P(X=k)=\binom{n}{k}p^k(1-p)^{n-k},\qquad
P(X=k)=e^{-\lambda}\frac{\lambda^k}{k!}
\]

\section{Unidad 6}
\[
f(x)=\frac{1}{\sigma\sqrt{2\pi}}\exp\Bigl(-\frac{(x-\mu)^2}{2\sigma^2}\Bigr),\qquad
Z=\frac{X-\mu}{\sigma}
\]

\section{Unidad 7}
\[
\bar{x}\pm z_{\alpha/2}\frac{\sigma}{\sqrt{n}},\qquad
\hat{p}\pm z_{\alpha/2}\sqrt{\frac{\hat{p}(1-\hat{p})}{n}}
\]

\section{Unidad 8}
\[
m_k=\frac{1}{n}\sum x_i^k,\qquad
\hat{\theta}_{\mathrm{MV}}=\arg\max_\theta \prod_{i=1}^n f(x_i\mid\theta)
\]

\end{document}
